\section{Computational Geometry}
\label{sec:geometry}

Every section so far has computed with numbers, polynomials, and functions. This
one computes with \emph{shapes}. A polygon is a list of corner points, a point
set is a scatter of coordinates, and the questions are the ones you would ask
with a ruler and a piece of string: how much area does this figure enclose, how
long is its boundary, where is its balance point, does a given point fall inside
it, and---given only a cloud of points---what is the smallest convex fence that
holds them all? These are the primitives of computational
geometry\index{computational geometry}, and Mathilda answers them for planar
figures built from \mcode{Polygon} and raw point lists.

The organizing principle is the one that runs through the whole system:
\emph{stay exact while the input is exact}. Give a polygon integer or rational
corners and its area comes back an exact integer or rational, its centroid an
exact pair of rationals, and its perimeter an exact sum of surds---never a
rounded decimal that has quietly discarded the last few digits. Under the hood
every coordinate is loaded into a GMP\index{GMP} rational and every predicate---
the sign of a cross product, the containment of a point---is decided by exact
integer comparison, so there are no tolerance thresholds and no near-miss
misclassifications. Only when you write a coordinate as an approximate real does
the computation switch to machine numbers, following the same numeric-contagion
rule as the rest of Mathilda.\index{exact arithmetic!in geometry}

We take the primitives in turn: first the two metric quantities of a single
polygon, its \emph{area} and \emph{perimeter} (\B{Area}, \B{Perimeter}); then its
\emph{centroid} (\B{RegionCentroid}); then the \emph{point-in-polygon} test
(\B{RegionMember}); and finally the \emph{convex hull} of a point set
(\B{ConvexHullRegion}), the one construction here that must invent its own output
shape rather than measure a given one.

\subsection{Area and perimeter of a polygon}
\label{ssec:geo-area}

The area of a polygon is computed by the \emph{shoelace formula}, and it is worth
seeing why it works before seeing it run. \B{Area} takes a \mcode{Polygon} and
returns the region it encloses:

\begin{codepairs}
\pair{geometry/area/1}{A four-by-three axis-aligned rectangle, stored for reuse.}
\pair{geometry/area/2}{\B{Area} returns the exact integer \(12\)---base times
height, as it must.}
\pair{geometry/area/3}{A triangle with a \emph{rational} corner: the area is the
exact rational \(1/4\), not \(0.25\). Nothing has been rounded.}
\pair{geometry/area/4}{A concave, five-sided polygon with a notch cut into its top
edge. The shoelace formula handles concavity with no special case: \(10\).}
\end{codepairs}

The rectangle and the concave pentagon make the same point from two directions.
The formula never asks whether the polygon is convex; it sums a signed
contribution from each edge, and the signs cancel exactly the double-counted and
re-entrant parts. And the rational triangle shows the exactness promise kept: an
integer-and-rational input yields a rational output, carried as \(1/4\) rather
than collapsed to a float.

\calloutnote{Theory}{The \emph{shoelace formula}\index{shoelace formula} gives the
area of a simple polygon with vertices \((x_1,y_1),\dots,(x_n,y_n)\) in order as
\[
  A \;=\; \frac12\left\lvert\sum_{i=1}^{n}\bigl(x_i\,y_{i+1}-x_{i+1}\,y_i\bigr)\right\rvert,
\]
indices taken cyclically so \((x_{n+1},y_{n+1})=(x_1,y_1)\). Each term is twice
the signed area of the triangle formed by the origin and one edge; summing over
the edges telescopes into the enclosed area by Green's
theorem\index{Green's theorem}, \(A=\tfrac12\oint(x\,dy-y\,dx)\). The name is the
mnemonic: writing the coordinates in two columns and multiplying along the
diagonals looks like lacing a shoe. Because every product is of two exact
coordinates, the sum---and hence the area---is computed with no rounding at all.}

\usagebox{Area}

The \B{Perimeter} is the sum of the edge lengths, the closing edge included. Here
the exactness promise takes a more interesting form, because a length is a square
root and most square roots are irrational:

\begin{codepairs}
\pair{geometry/perimeter/1}{A unit right triangle: two legs of length \(1\) and a
hypotenuse of \(\sqrt2\), returned as the exact surd \(2+\sqrt2\).}
\pair{geometry/perimeter/2}{Shrink one leg to \(1/2\) and the hypotenuse becomes
\(\sqrt{5}/2\); the perimeter is the exact \(3/2+\tfrac12\sqrt5\).}
\pair{geometry/perimeter/3}{Write even one corner as a real---here a \(3\)-\(4\)-\(5\)
triangle with real coordinates---and the whole computation goes to machine
numbers: \(12.0\).}
\end{codepairs}

Where \B{Area} could keep everything rational, \B{Perimeter} cannot: the length of
an edge \((\Delta x,\Delta y)\) is \(\sqrt{\Delta x^2+\Delta y^2}\), and Mathilda
builds exactly that expression---a sum of \mcode{Sqrt} terms---and hands it to the
evaluator, which canonicalizes the radicals (\(\sqrt8\to2\sqrt2\),
\(\sqrt{5/4}\to\tfrac12\sqrt5\)) and collects the result. So the perimeter of the
first triangle is the symbolic \(2+\sqrt2\), not \(2.414\dots\); you get digits
only when you ask, with \B{N}. The moment a coordinate is inexact, though, there
is no exactness left to preserve, and the third example drops straight to a
machine \(12.0\).\index{exact arithmetic!surds}

\calloutnote{Under the hood}{Both heads live in \texttt{src/geometry.c}. Exact
coordinates are parsed into GMP rationals (\texttt{mpq\_t}); \B{Area} accumulates
the shoelace sum in rational arithmetic and normalizes the result to an
\mcode{Integer} when the denominator is \(1\). \B{Perimeter} instead assembles a
\mcode{Plus} of \mcode{Sqrt}\texttt{[dx\^{}2 + dy\^{}2]} terms and lets the
ordinary evaluator do the radical simplification, which is why its output looks
exactly like arithmetic you could have typed by hand. A single real coordinate
flips an \texttt{exact} flag and the same routines run on \texttt{double}s
instead. Consecutive repeated vertices---including the wrap-around pair---are
collapsed first, so a degenerate edge cannot smuggle a spurious there-and-back
length into the sum; a figure that collapses to fewer than three distinct corners
returns \mcode{Undefined} rather than a confident \(0\).}

\subsection{The centroid}
\label{ssec:geo-centroid}

The \emph{centroid} of a region is its center of mass, assuming uniform density---
the point on which a cardboard cutout of the polygon would balance.
\B{RegionCentroid} returns it as a coordinate pair, exact when the polygon is:

\begin{codepairs}
\pair{geometry/centroid/1}{The centroid of the unit right triangle sits at
\(\{1/3,1/3\}\)---one third of the way in from each leg, the classical result that
the centroid divides each median \(2\!:\!1\).}
\pair{geometry/centroid/2}{For the concave pentagon the balance point is
\(\{2,7/5\}\): pulled below the center by the notch removed from the top, and
exact to the last digit.}
\end{codepairs}

The triangle answer is one you can check against a theorem: the centroid of a
triangle is the average of its three vertices,
\(\tfrac13(\{0,0\}+\{1,0\}+\{0,1\})=\{1/3,1/3\}\). For a general polygon the
centroid is \emph{not} the average of the vertices---that would weight a cluster
of nearby corners too heavily---but an area-weighted integral, and the concave
pentagon shows the difference: its answer reflects where the mass actually is,
notch and all.

\calloutnote{Theory}{For a simple polygon the centroid \((C_x,C_y)\) is the
area-weighted average
\[
  C_x=\frac{1}{6A}\sum_{i=1}^{n}(x_i+x_{i+1})\,(x_i y_{i+1}-x_{i+1}y_i),
  \qquad
  C_y=\frac{1}{6A}\sum_{i=1}^{n}(y_i+y_{i+1})\,(x_i y_{i+1}-x_{i+1}y_i),
\]
with \(A\) the signed shoelace area and the same cyclic indexing. The cross-product
weights \(x_i y_{i+1}-x_{i+1}y_i\) are exactly the shoelace terms of the previous
subsection, so the centroid reuses that computation; dividing two exact rational
sums keeps the result rational.\index{centroid!of a polygon} A zero-area polygon
has no area to weight by, so \B{RegionCentroid} declines rather than dividing by
zero.}

\subsection{Point-in-polygon: the membership test}
\label{ssec:geo-member}

Given a polygon and a point, is the point inside it? This innocent question is a
workhorse of graphics, geographic information systems, and collision detection,
and its exact answer on the boundary is where naive implementations go wrong.
\B{RegionMember} returns \mcode{True} when the point lies inside \emph{or on} the
polygon, \mcode{False} otherwise:

\begin{codepairs}
\pair{geometry/member/1}{A \(2\times2\) square, stored for the queries below.}
\pair{geometry/member/2}{A point squarely in the interior: \mcode{True}.}
\pair{geometry/member/3}{A point \emph{on the right edge}, \(\{2,1\}\):
\mcode{True}. Boundary points count as members.}
\pair{geometry/member/4}{A point outside the square: \mcode{False}.}
\pair{geometry/member/5}{An exact \emph{rational} query point, \(\{1,1\}/2\), inside
a triangle: decided exactly, \mcode{True}.}
\pair{geometry/member/6}{The subtle case---a point in the \emph{notch} of the
concave pentagon. It sits within the polygon's bounding box but outside the figure
itself, and the crossing test correctly returns \mcode{False}.}
\end{codepairs}

Two of these repay attention. The boundary point \(\{2,1\}\) returns \mcode{True}
because the test is deliberately inclusive, and---crucially---because the input is
exact, that boundary decision is exact: the point lies on the edge if and only if
a cross product is \emph{exactly} zero, a question GMP integers answer with no
ambiguity. A floating-point implementation would have to guess at a tolerance and
could fall either way. And the notch case shows why a bounding-box check is not
enough: the concave figure requires a genuine crossing test, which handles
concavity as transparently as the shoelace formula did.

\calloutnote{Theory}{The membership test is the \emph{ray-casting}\index{ray casting}
algorithm, a direct application of the Jordan curve
theorem\index{Jordan curve theorem}: a simple closed curve divides the plane into
an inside and an outside, and a ray from any point crosses the boundary an odd
number of times if the point is inside, an even number if outside. Mathilda casts a
horizontal ray and counts the edges it crosses, using the sign of the edge--point
cross product to decide each crossing consistently; a separate exact check catches
the point lying directly on an edge (cross product zero, and inside the edge's
bounding box) and reports membership. On exact input every one of these decisions
is an exact integer comparison---the code path in \texttt{src/geometry.c} runs the
crossing count over \texttt{mpq\_t} coordinates---so there is no epsilon and no
boundary case that rounds the wrong way.}

\subsection{The convex hull}
\label{ssec:geo-hull}

The constructions so far measured a polygon you supplied. \B{ConvexHullRegion}
does the opposite: given only a scatter of points, it \emph{builds} the smallest
convex region containing them all---the shape a rubber band would take if
stretched around the whole set and released. It is the most algorithmically
substantial primitive in this section, and its result depends on the dimension of
the points' arrangement:

\begin{codepairs}
\pair{geometry/hull/1}{Six points: the four corners of a square, one point on the
bottom edge, and one dead center.}
\pair{geometry/hull/2}{The hull is the square. The mid-edge point \(\{1,0\}\) is
collinear and the center point \(\{1,1\}\) is interior, so both are dropped---only
the four extreme corners survive, listed counterclockwise from the lexicographic
minimum.}
\pair{geometry/hull/3}{Feeding the hull straight into \B{Area} gives \(4\): the
primitives compose, and the whole pipeline stays exact.}
\pair{geometry/hull/4}{The four corners of a square with its center added again
reduce to the square---the interior point never reaches the boundary.}
\pair{geometry/hull/5}{When the points are all \emph{collinear} there is no
two-dimensional hull, so the result degrades to the \mcode{Line} through the two
extremes.}
\pair{geometry/hull/6}{And a single point is its own hull, returned as a
\mcode{Point}.}
\end{codepairs}

The first example is the whole idea in miniature. Of six input points only four
are \emph{vertices} of the hull; the collinear and interior points contribute
nothing to the convex boundary and are discarded. That the vertices come back in
counterclockwise order starting from the lexicographically smallest point is not
incidental---it is a canonical ordering, so the same point set always yields the
same \mcode{Polygon}, and two hulls can be compared directly. The degenerate
returns matter too: rather than forcing a two-dimensional answer onto
one-dimensional data, the routine hands back a \mcode{Line} or a \mcode{Point},
telling you the true dimension of the arrangement.

\calloutnote{Under the hood}{\B{ConvexHullRegion} uses \emph{Andrew's monotone
chain}\index{convex hull!monotone chain}\index{Andrew's monotone chain}, a clean
variant of the Graham scan. The points are first sorted lexicographically by
\(x\) then \(y\); duplicates, now adjacent, are removed. The algorithm then sweeps
left to right building the \emph{lower} hull and right to left building the
\emph{upper} hull, at each step popping any vertex that would make a non-left
turn---detected by the sign of the cross product of the last two edges. A
cross product of zero means three collinear points, and the middle one is popped,
which is exactly why interior and collinear-middle points never appear in the
result. The cost is \(O(n\log n)\), dominated by the sort; the hull-building sweeps
are linear. As everywhere in \texttt{src/geometry.c}, exact input runs the turn
predicate on \texttt{mpq\_t} cross products, so collinearity is decided
exactly---no nearly-collinear point is kept or dropped by a rounding accident.}

\usagebox{ConvexHullRegion}

\calloutnote{Pitfall}{These heads assume a \emph{simple} polygon---one whose edges
do not cross. Hand \B{Area} a self-intersecting vertex list and it returns the
signed shoelace value under the even--odd rule, which is not the area Wolfram
Language reports for the enclosed region; self-intersection is not detected at
runtime. They are also strictly planar: a non-2D point list, symbolic
coordinates, or a \mcode{Polygon} with holes leaves the expression unevaluated
rather than guessing. And a coordinate too large to represent as a machine double,
queried on the machine path, is declined outright---an overflow to infinity would
turn every cross product into a \mcode{NaN} that silently reads as ``collinear,''
so the honest answer is to decline. Keep the coordinates exact and that trap never
arises.}

\subsection*{Where this connects}

Computational geometry is where Mathilda's exact-arithmetic core meets the plane.
Its bedrock is the rational arithmetic of \S\ref{sec:arithmetic}: every area,
centroid, and membership decision is ultimately a comparison of GMP rationals, and
the surds returned by \B{Perimeter} are canonicalized by the same radical
machinery that serves the algebra of \S\ref{sec:algebra}. The cross product that
decides a turn in the convex hull and a crossing in the membership test is the
same \(2\times2\) determinant that the linear algebra of \S\ref{sec:linalg}
computes in the large, and the shoelace area is a discrete line integral---a
foothold on the vector calculus of the numerical chapter. Once a region is in
hand, the graphics chapter (Chapter~\ref{ch:graphics}) draws it: a \mcode{Polygon}
or the \mcode{Line} and \mcode{Point} that \B{ConvexHullRegion} can return are
themselves graphics primitives, so a hull can be computed exactly and then shown.
For the exhaustive options of any function named here, its reference page is one
\texttt{?Name} away at the prompt.
